\documentclass{article}

\def\CourseCode{ECEN-662}
\def\CourseName{Estimation \& Detection Theory}
\def\Instructor{Dr. Tie Liu}
<<<<<<< HEAD
\def\Semester{Spring 2026}
\def\Student{Brody Jordan}
\def\University{Texas A\&M University}
\def\Cover{figures/_cover.pdf}
=======
\def\Semester{Spring 2026}
>>>>>>> 249e4c317b2a4234368ca0c3c75e714883e8075a

\usepackage[
    AssignmentName={Homework 1},
    DueDate={January 27, 2026 at 11:59 PM}
]{../../../styles/assignment}

\begin{document}
\MakeAssignmentTitle

% ------ PROBLEM 1 ------ %

% \[
% \begin{cases}
%   1 - \theta, x = 0 \\
%   \frac{\theta}{2}, x = \pm 1
% \end{cases}
% \] 
% write as exponential family, show family is full-ranked, shows that sufficient statistic is also complete...

% $\mathbb{P} [\phi(x) = 0] < 1$ expectation is zero, but random variable itself is not always zero

\begin{homeworkProblem}

    Let $X \sim p_\theta(x)$, where
    \[
        p_\theta(x) = \left(\frac{\theta}{2}\right)^{|x|} (1 - \theta)^{1 - |x|}, \quad x \in \{-1, 0, 1\} \quad \text{and} \quad \theta \in (0, 1)
    \]
    \begin{enumerate}[label=\roman*), topsep=5pt, itemsep=5pt]
        \item Does $p_\theta(x)$ belong to the exponential family?
        \item Is $|X|$ a complete sufficient statistic?
        \item Is $X$ complete? \newline
    \end{enumerate}

    \solution

    \textbf{Part (i):} This pmf can be rewritten into a different form
    \begin{align*}
        p_\theta(x) & = \exp \left( \log \left( \left(\frac{\theta}{2}\right)^{|x|} (1 - \theta)^{1 - |x|} \right) \right)                            \\
                    & = \exp \left( |x| \log \frac{\theta}{2} + \log (1 - \theta) - |x| \log (1 - \theta) \right)                                     \\
                    & = \exp \left( \underbrace{|x|}_{T(x)} \underbrace{\log \frac{\theta}{2(1 - \theta)}}_{\eta(\theta)} + \log (1 - \theta) \right)
    \end{align*}
    Since the model $p_\theta(x)$ can be expressed in the form $f_\theta(x) = h(x) \exp(\eta^t T - A(\theta))$, it belongs to the exponential family. \\

    \textbf{Part (ii):} $T(X) = |X|$ is a complete sufficient statistic, since $\E_\theta[\phi(T)] = 0 ,\ \forall \theta \in (0, 1)$.  \\

    \textbf{Part (iii):} $X$ is not complete since $\E_\theta[\phi(T)] = 0, \ \forall \theta \in (0, 1)$ does not imply that $ \mathbb{P}_\theta(\phi(T) = 0) = 1, \ \forall \theta \in (0, 1)$.

\end{homeworkProblem}

\pagebreak

% ------ PROBLEM 2 ------ %
\begin{homeworkProblem}

    Let $X = (X_1, \dots, X_n)$, where $X_i$'s are i.i.d. $\mathcal{N}(\theta, \theta)$, and $\theta > 0$ is an unknown parameter. Determine a complete sufficient statistic for $\theta$ and carefully justify your answer. \\

    \solution

    The pdf of $\mathcal{N}(\theta, \theta)$ is given by
    \begin{align*}
        f_\theta(x_i) & = \frac{1}{\sqrt{2 \pi \theta}} \exp \left( -\frac{(x_i - \theta)^2}{2 \theta} \right)                                  \\
                      & = \frac{1}{\sqrt{2 \pi \theta}} \exp \left( -\frac{x_i^2}{2 \theta} + x_i - \frac{\theta}{2} \right)                    \\
                      & = \frac{1}{\sqrt{\pi}} \exp \left( -\frac{x_i^2}{2\theta} + x_{i} + \frac{\theta}{2} -\frac{1}{2} \log(2\theta) \right)
    \end{align*}
    thus $\eta(\theta) = \left( 1, -\frac{1}{2\theta}\right)$ and $T(x_i) = (x_i, x_i^2)$, which is a complete sufficient statistic for $\theta$.

\end{homeworkProblem}

\pagebreak

% ------ PROBLEM 3 ------ %
\begin{homeworkProblem}

    Let $X = (X_1, \dots, X_n)$, where $X_i$'s are i.i.d. $\mathcal{N}(\theta, \theta^2)$, and $\theta > 0$ is an unknown parameter.
    \begin{enumerate}[label=\roman*), topsep=5pt, itemsep=5pt]
        \item For $n = 1$, show that $T = X$ is a minimum sufficient statistic but is not a complete sufficient statistic.
        \item Show that the sufficient statistic
              \[
                  T = \left(\sum_{i=1}^{n} X_i, \sum_{i=1}^{n} X_{i}^2 \right)
              \]
              is not complete for any $n \geq 2$. \newline
    \end{enumerate}

    \solution

    \textbf{Part (i):} Here
    \begin{align*}
        f_\theta(x_i) & = \frac{1}{\sqrt{2 \pi \theta^2}} \exp \left( -\frac{(x_i - \theta)^2}{2 \theta^2} \right)          \\
                      & = \frac{1}{\sqrt{2 \pi \theta^2}} \exp \left( -\frac{x_i^2}{2 \theta^2} + x_i - \frac{1}{2} \right)
    \end{align*}
    so $\eta(\theta) = \left( 1, -\frac{1}{2\theta^2} \right)$ and $T(x_i) = (x_i, x_i^2)$, the dimensionality difference implies that $T$ is not a complete sufficient statistic. \\

    \textbf{Part (ii):} When $T = \left(\sum_{i=1}^{n} X_i, \sum_{i=1}^{n} X_{i}^2 \right)$, there exists a $\phi(T)$ such that $\E_\theta[\phi(T)] = 0$ but $\mathbb{P}_\theta(\phi(T)=0) = 1, \ \forall \theta \in \Theta$, which means $T$ is not a complete statistic when $n \geq 2$. \\

\end{homeworkProblem}

\pagebreak

% ------ PROBLEM 4 ------ %
\begin{homeworkProblem}

    Prove the following law of total variance:
    \[
        \Var[Y] = \E[\Var[Y|Z]] + \Var[\E[Y|Z]]
    \]
    for any real-valued random variable $Y$. \\

    \solution

    Starting from the variance definition, and applying expectation and variance properties:
    \begin{align*}
        \Var[Y] & = \E[Y^2] - E[Y]^2                              \\
                & = \E[\E[Y^2|Z]] - \E[\E[Y|Z]]^2                 \\
                & = \E[\Var[Y|Z] + \E[Y|Z]^2] - \E[\E[Y|Z]]^2     \\
                & = \E[\Var[Y|Z]] + \E[\E[Y|Z]^2] - \E[\E[Y|Z]]^2 \\
                & = \E[\Var[Y|Z]] + \Var[\E[Y|Z]]
    \end{align*}
    Which is the law of total variance.

\end{homeworkProblem}

\end{document}




% % Problem 2
% Let $X = (X_1, \dots, X_n)$, where $X_i$'s are i.i.d. $\mathcal{N}(\theta, \theta)$, and $\theta > 0$ is an unknown parameter. Determine a complete sufficient statistic for $\theta$ and carefully justify your answer.
% \solution 
% start by checking whether this is an expoential family or not...

% % Problem 3 
% Let $X = (X_1, \dots, X_n)$, where $X_i$'s are i.i.d. $\mathcal{N}(\theta, \theta^2)$, and $\theta > 0$ is an unknown parameter.
% \begin{enumerate}
%   \item For $n=1$, show that $T=X$ is a minimum sufficient statistic but is not a complete sufficient statistic.
%   \item Show that the sufficient statistic
%     \[
%       T = \left( \sum_{i=1}^{n} X_i, \sum_{i=1}^{n} X_{i}^2 \right)
%     \]
%     is not complete for any $n \geq 2$.
% \end{enumerate}
% able to compute $\E_\theta[(2x_i)^2]$
% $\E_\theta[\sum X_i^2$